% !TeX root = ../../main.tex
\subsection{英作文の問題について}
    英作文は北海道公立高校入試において、各大問の\fitblankbf[]{最後の問題}に用意されている。
    配点は英作文がメインの大問4で\fitblankbf[]{12}点。試験全体を通して大問4を含めて\fitblankbf[]{40}点分配点されている。
    そのうちのほとんどは、「\fitblankbf[]{主語と動詞を含む}1文で答えなさい」という形式で出題される。

    英作文は基本的に\fitblankbf[]{減点方式}で採点される。したがって、\uwave{単語や文法の間違いをなるべく避ける}
    ことが得点の鍵である。その具体的な方法はのちに示すこととする。

    大問4は3題あり、そのうち2題は\fitblankbf[]{空欄補充}、最後の1題が\fitblankbf[]{自由英作文}である。

\subsection{北海道公立高校入試の英作文における一般的な対策}
    リスニング同様、毎年、配点も問題構成も変わらない。仮に変わったとしても、全員が同じ問題を解くことになり、
    それを予想できた人間はいないはずであるから、次に示す対策をしない理由はない。

    \begin{enumerate}
        \item \fitblankbf[]{問題構成}をしっかりと理解する。
        \item \fitblankbf[]{24}語以上の英語で書かなければ、\fitblankbf[]{採点対象}にならない。
        \item 慣れてさえしまえば、長文読解に比べて、難易度的にも、時間的にも\fitblankbf[]{得点しやすい}。
    \end{enumerate}

\subsection{英作文における一般的な対策}
    まず、大前提として、英作文は\uwave{自分が表現できる英語で書く}ことが極めて重要である。
    具体的には、\uwave{書く内容が嘘}でも良いし、\uwave{幼稚な内容}だって構わない。
    なぜなら、採点の基準は\uwave{語数}と\uwave{単語、文法の正確さ}であり、内容が概ね本題に沿っていれば
    \uwave{内容の正しさ}は無関係であるからである。

    24語程度なら

    特に、英作文をする上で必ず押さえておきたいポイントを次にまとめる。
\newpage

        \begin{techniquebox}{日々の学習で基礎語彙力、基礎文法力を身につけよう!}

        \end{techniquebox}
\vspace{1em}

        \begin{techniquebox}{戦える武器を持とう!}

        \end{techniquebox}
\vspace{1em}

        \begin{techniquebox}{日本語を工夫しよう!}

        \end{techniquebox}
\vspace{1em}
